\documentclass[10pt,a4paper]{article}
\usepackage[utf8]{inputenc}
%\usepackage[german]{babel}
%\usepackage[cyrillic]{babel}
\usepackage[T1]{fontenc}
\usepackage{amsmath}
\usepackage{amsfonts}
\usepackage{amssymb}
\usepackage{amsthm}
%\usepackage{stix-type1}%\turnediota
\theoremstyle{definition}
\newtheorem{definition}{Definition}
\newtheorem{lemma}{Lemma}
\newtheorem{theorem}{Theorem}
\usepackage{etoolbox}

%\newcommand*{\@Dashv}[2]{%
%  \reflectbox{$\m@th#1#2$}%
%}
\newcommand*{\separation}[3]{% 1)binder, 2)set, 3)condition
  \{#1\in #2: #3\}
}

% https://www.freecodecamp.org/news/simple-regex-tricks-for-beginners-3acb3fa257cb/
%  \\mathcal\{T\} |---> \topology
%  \\mathcal\{([^\}]*)\}   |--->  $1
\newcommand{\complex}{\mathbb{C}}
\newcommand{\real}{\mathbb{R}}
\newcommand{\power}{\mathcal{P}}
\newcommand{\borel}{\mathcal{B}}
\newcommand{\interior}{\mbox{int}}
\newcommand{\closure}{\mbox{cl}}
\newcommand{\nghbset}{\mathcal{N}}
\newcommand{\semialgebra}{\mathcal{S}}
\newcommand{\topology}{\mathcal{T}}
\newcommand{\union}{\bigcup}
\newcommand{\functional}{\mbox{ functional}}
\newcommand{\injective}{\mbox{ }}
\newcommand{\total}{\mbox{ total on }}
\newcommand{\surjective}{\mbox{ surjective on}}
\newcommand{\countable}{\mbox{ countable}}
\newcommand{\converges}{\longrightarrow}
\newcommand{\AND}{\ \&\ }
\newcommand{\IMPLIES}{\ \Rightarrow\ }
\newcommand{\Proof}{\noindent \textbf{Proof.}\\}
\newcommand{\Done}{\noindent \textbf{Done.}\\}
% local commands
%\preto\section{\undef\blipo} %etoolbox
\preto\subsection{\undef\blipo} %etoolbox
% https://tex.stackexchange.com/questions/105684/limit-macro-newcommand-scope-to-section
\begin{document}
\section{Logic}
Partially based on WILLIAM M. FARMER \& JOSHUA D. GUTTMAN's article\\
``A Set Theory with Support for Partial Functions''.\\
For simplicity of understanding by the reader, the notation from the article is used here.\\
\begin{definition} Value of a function.
$$f(x) \simeq \big(\iota y: \langle x,y\rangle\in f\big)$$
\end{definition}

\begin{theorem} If $f$ is univocal, then
$$\big(f(x)\simeq y\big)\Leftrightarrow \big(\langle x,y\rangle\in f\big)$$
\end{theorem}
\Proof
($\Leftarrow$) (Obvious part) Here we use axiom PFOL$10'$ which is equaivalent to PFOL10:
$$
(\iota x:\varphi)\downarrow
\quad\overset{\mbox{PFOL10}'}{\Longleftrightarrow}\quad
(\exists!x:\varphi)
$$
So the proof is the following:
$$
\big(\langle x,y\rangle\in f\big)
\quad\overset{\exists\mbox{-intro}}{\Longrightarrow}\quad
\big(\exists w: \langle x,w\rangle\in f\big)
$$
$$
f\mbox{ univocal}
\quad\overset{\mbox{def}}{\Longleftrightarrow}\quad
\big(! w: \langle x,w\rangle\in f\big)
$$
$$
\big(\exists w: \langle x,w\rangle\in f\big)
\quad\&\quad
\big(! w: \langle x,w\rangle\in f\big)
\quad\overset{\exists!\mbox{-intro}}{\Longrightarrow}\quad
\big(\exists! w: \langle x,w\rangle\in f\big)
$$
$$
\big(\exists! w: \langle x,w\rangle\in f\big)
\quad\overset{\mbox{PFOL10}'}{\Longrightarrow}\quad
\big(\iota w: \langle x,w\rangle\in f\big)\downarrow
$$
$$
\big(\iota w: \langle x,w\rangle\in f\big)\downarrow 
\quad\overset{\mbox{PFOL11}}{\Longrightarrow}\quad
\langle x,(\iota w: \langle x,w\rangle\in f)\rangle\in f
$$
$$
\langle x,(\iota w: \langle x,w\rangle\in f)\rangle\in f
\quad\&\quad
\big(\langle x,y\rangle\in f\big)
\quad\overset{f \mbox{ univocal}}{\Longrightarrow}\quad
(\iota w: \langle x,w\rangle\in f) \simeq y
$$
$$
f(x) \simeq y
$$
%\big(! w: \langle x,w\rangle\in f\big)
%$$
%\big(\iota w: \langle x,w\rangle\in f\big)\downarrow

%$$
\noindent($\Rightarrow$)
$$(y\downarrow)\quad\mbox{(by PFOL8)}$$
$$
(y\downarrow)\quad\&\quad((\iota w: \langle x,w\rangle\in f)\simeq y) \quad\overset{\mbox{PFOL7}}{\Longrightarrow}\quad
(\iota w: \langle x,w\rangle\in f)\downarrow $$
$$
\big(\iota w: \langle x,w\rangle\in f\big)\downarrow 
\quad\overset{\mbox{PFOL11}}{\Longrightarrow}\quad
\langle x,(\iota w: \langle x,w\rangle\in f)\rangle\in f
$$
$$
\langle x,(\iota w: \langle x,w\rangle\in f)\rangle\in f
\quad\&\quad
\big((\iota w: \langle x,w\rangle\in f)\simeq y\big)
\quad\overset{\mbox{PFOL7}}{\Longrightarrow}\quad
\big(\langle x, y\rangle\in f\big)
$$
\\
\Done

\end{document}\\
